\documentclass[11pt, a4paper]{article}

\usepackage[utf8]{inputenc}
\usepackage[T1]{fontenc}
\usepackage[spanish]{babel}
\usepackage[margin=2cm, headheight=25pt]{geometry}
\usepackage{fancyhdr}
\usepackage{xcolor}
\usepackage{colortbl}
\usepackage{tabularx}
\usepackage{booktabs}
\usepackage{tcolorbox}
\usepackage{amssymb}
\usepackage{enumitem} % Requerido para modificar espacios en listas [itemsep=...]
\usepackage{microtype} % Optimiza la justificación y evita Overfull \hbox

\definecolor{ucbblue}{RGB}{0, 51, 102}

\tcbset{
    caja_instruccion/.style={colback=blue!5, colframe=ucbblue, fonttitle=\bfseries, boxrule=1pt, arc=3pt}
}

\pagestyle{fancy}
\fancyhf{}
\lhead{\textcolor{ucbblue}{\textbf{IMT-121 | Instrumento de Defensa}}}
\rhead{\textbf{Hito 1 -- PFI}}
\cfoot{Página \thepage}
\renewcommand{\headrulewidth}{0.4pt}

\begin{document}

\begin{center}
    {\LARGE \textbf{Hoja de Evaluación y Defensa: Hito 1 (PFI)}} \\[0.2cm]
    {\large Cierre de Evaluación y Diagnóstico de Competencias (EC1)}
\end{center}

\begin{tcolorbox}[caja_instruccion, title=Regla de Equidad (Fairness Rule)]
Diferenciar estrictamente: \textbf{[S] Sumativo} (comunicado explícitamente a los estudiantes y con impacto en la nota) vs. \textbf{[D] Diagnóstico} (evidencia deseable para medir el estado del curso, pero sin impacto punitivo retroactivo).
\end{tcolorbox}

\vspace{0.2cm}
\noindent\textbf{Equipo N$^\circ$:} \rule{3cm}{0.4pt} \hspace{1cm} \textbf{Track:} $\square$ Mecatrónica \quad $\square$ Biomédica \\[0.2cm]
\textbf{Integrantes:} 1) \rule{4.2cm}{0.4pt} 2) \rule{4.2cm}{0.4pt} 3) \rule{4.2cm}{0.4pt}

\section*{1. Banco de Preguntas Estandarizadas (Rotar entre integrantes)}
\textit{Instrucción: El estudiante debe responder en 1-2 minutos en la pizarra/pantalla.}
\begin{itemize}[itemsep=2pt]
    \item \textbf{Circuito a Matriz:} ``Muestra físicamente de dónde proviene esta fila/columna de la matriz.''
    \item \textbf{Sensibilidad:} ``Si esta resistencia $R_k$ cambia, ¿qué celdas exactas de la matriz cambian y por qué?''
    \item \textbf{Código:} ``¿Cómo garantizas que el orden de tus variables en Python coincide con la matriz planteada?''
    \item \textbf{Simulación:} ``¿Qué parámetro en la salida \texttt{.op} de LTspice corresponde a esta incógnita analítica?''
    \item \textbf{Física:} ``Si la simulación y el hardware difieren más del 5\%, ¿cuál es el primer instrumento que auditarías?''
\end{itemize}

\section*{2. Rejilla de Observación y Captura de Evidencia}
\renewcommand{\arraystretch}{1.5}
\noindent
\begin{tabularx}{\textwidth}{|c|X|c|c|c|c|}
    \hline
    \rowcolor{ucbblue!10}
    \textbf{Tipo} & \textbf{Criterio Evaluado} & \textbf{Logrado} & \textbf{Parcial} & \textbf{No Evi.} & \textbf{N/A} \\
    \hline
    [S] & \textbf{Física:} Interpreta el circuito y plantea ecuaciones base. & $\square$ & $\square$ & $\square$ & $\square$ \\
    \hline
    [S] & \textbf{Matriz:} Ensambla $\mathbf{A}\mathbf{x}=\mathbf{b}$ correctamente (Mallas/Nodos). & $\square$ & $\square$ & $\square$ & $\square$ \\
    \hline
    [S] & \textbf{Simulación:} Interpreta salidas \texttt{.op} en LTspice. & $\square$ & $\square$ & $\square$ & $\square$ \\
    \hline
    [D] & \textbf{Software:} Domina \texttt{numpy.linalg.solve} con autonomía. & $\square$ & $\square$ & $\square$ & $\square$ \\
    \hline
    [D] & \textbf{Hardware:} Validó físicamente con $\varepsilon_r \le 5.0\%$. & $\square$ & $\square$ & $\square$ & $\square$ \\
    \hline
    [S] & \textbf{Defensa:} Explica discrepancias y defiende el modelo. & $\square$ & $\square$ & $\square$ & $\square$ \\
    \hline
\end{tabularx}

\section*{3. Registro de Debilidades Críticas (Carry-Forward)}
\textit{Marque las áreas que requieren refuerzo embebido en la Unidad 2:}

\vspace{0.3cm}
\noindent
\renewcommand{\arraystretch}{1.3}
\begin{tabular}{@{}ll@{}}
    $\square$ Falla en Ley de Ohm/Kirchhoff & $\square$ Errores de signos en Matriz \\
    $\square$ Confusión Mallas/Nodos & $\square$ Copia de código sin entender \\
    $\square$ No correlaciona LTspice con el modelo & $\square$ Hardware inestable
\end{tabular}

\vspace{0.6cm}
\noindent\textbf{Decisión del Hito 1:} $\square$ Aprobado \quad $\square$ Requiere Ajustes \quad \textbf{Puntaje Asignado:} \rule{2cm}{0.4pt} / 100

\end{document}
